\documentclass[11pt]{article}
\usepackage[a4paper,margin=18mm]{geometry}
\usepackage{fontspec}
\setmainfont{Latin Modern Roman}
\usepackage{amsmath,amssymb,amsthm,amscd,graphicx,color}
\usepackage{xurl}
\usepackage[unicode,hidelinks]{hyperref}
\allowdisplaybreaks
\setlength\emergencystretch{3em}

\makeatletter
\newtheorem{theorem}{Theorem}
\newtheorem*{theorem*}{Theorem}
\newtheorem{lemma}{Lemma}
\newtheorem*{lemma*}{Lemma}
\newtheorem{corollary}{Corollary}
\newtheorem*{corollary*}{Corollary}
\newtheorem{proposition}{Proposition}
\newtheorem*{proposition*}{Proposition}
\newtheorem{conjecture}{Conjecture}
\newtheorem*{conjecture*}{Conjecture}
{\theoremstyle{definition} \newtheorem{definition}{Definition}
\newtheorem*{definition*}{Definition}
\newtheorem{example}{Example}
\newtheorem*{example*}{Example}
\newtheorem{remark}{Remark}
\newtheorem*{remark*}{Remark}
\newtheorem{note}{Note}
\newtheorem*{note*}{Note}
}
\makeatother
\usepackage[all]{xy}
\numberwithin{theorem}{section}
\numberwithin{proposition}{section}
\numberwithin{lemma}{section}
\numberwithin{corollary}{section}
\numberwithin{definition}{section}
\numberwithin{example}{section}
\numberwithin{remark}{section}
\numberwithin{note}{section}

\DeclareMathOperator{\Ker}{Ker}
\DeclareMathOperator{\End}{End}

\numberwithin{equation}{section}

\begin{document}
\begin{center}\Large Orthogonal and even-dimensional symplectic twisted Yangians over C[hbar] as the specified graded limits of twisted quantum loop algebras\end{center}
\noindent\textbf{Stable ID:} 2008\par
\noindent\textbf{Authors:} Patrick Conner; Nicolas Guay\par
\noindent\textbf{Original work:} From Twisted Quantum Loop Algebras to Twisted Yangians\par
\noindent\textbf{Version:} Official SIGMA 11 (2015) article 040, DOI 10.3842/SIGMA.2015.040; journal PDF and native TeX acquired 2026-10-01, exact bytes pinned by SHA256\par
\noindent\textbf{Selected task scope:} Full Theorem3.2 for g\_N=\allowbreak{}o\_N or sp\_N over C (N positive, even in the symplectic case). The twisted Yangian of Definition3.1 over C[hbar] is isomorphic via s\_ij\textasciicircum{}(m+\allowbreak{}1)->S\_ij\textasciicircum{}(1,m) modulo K\_(m+\allowbreak{}1)\textasciicircum{}tw to the explicitly defined graded limit in Section3; hbar is the graded class of q-q inverse. No claim for an unfiltered or generic-q algebra.\par
\noindent\textbf{Difficulty:} H2: The proof constructs a filtration adapted simultaneously to q approaching1 and the loop parameter approaching1, then identifies the leading twisted RTT generators. Its supporting untwisted degeneration proves injectivity by tensor evaluation and differential separation of PBW monomials, while a spanning argument establishes surjectivity of the twisted associated graded map.\par
\noindent\textbf{Editorial boundary note (not author text):} The complete original twisted graded-limit isomorphism Theorem 3.2 and all author twisted and supporting untwisted degeneration proof are retained. Both o\_N and sp\_N cases keep their original RTT/integral-form conventions, with even N and the q-skew B matrix in the symplectic case. The full two-parameter filtration, evaluation/differential separation of untwisted PBW monomials, twisted compatibility and leading-generator spanning/surjectivity calculation are supplied. The source maps s\_ij\textasciicircum{}(m+\allowbreak{}1) to the class of S\_ij\textasciicircum{}(1,m) in K\_m\textasciicircum{}tw/K\_(m+\allowbreak{}1)\textasciicircum{}tw; hbar is the class of q-q inverse. Prior RTT embeddings/specialization are cited established background; the corresponding symplectic spanning computation is the bounded author-described adaptation. No naive q=\allowbreak{}1, unfiltered/generic-q, exceptional-type or super-Yangian extension is selected.\par
\noindent\textbf{Source:} \url{https://sigma-journal.com/2015/040/}\quad DOI: \url{https://doi.org/10.3842/SIGMA.2015.040}\par
\noindent\textbf{License and attribution:} Copyright retained by the named authors. Published by SIGMA under CC BY 4.0: \url{https://creativecommons.org/licenses/by/4.0/}. Source: \url{https://sigma-journal.com/about.html}. Adaptation consists of standalone excerpt formatting, cover and numbering restoration; original author language and mathematical text are retained.\par
\medskip\hrule\medskip\noindent\textbf{Author text begins below.}\par
\setcounter{section}{1}
\setcounter{subsection}{0}
\setcounter{subsubsection}{0}
\setcounter{equation}{0}
% Original source lines 142--977
\section[Quantum loop algebra and Yangian for $\mathfrak{gl}_N$]{Quantum loop algebra and Yangian for $\boldsymbol{\mathfrak{gl}_N}$}
%\label{glnsec}

We start by recalling the def\/inition of the Yangian and of the quantum loop algebra of $\mathfrak{gl}_N$ in
the RTT presentation~\cite{FRT}.
We view $\hbar$ as a~formal parameter.

\begin{definition}
\label{Ygln}
Let $R(u) = 1 \otimes 1 - \hbar u^{-1} \sum\limits_{i,j=1}^N E_{ij} \otimes E_{ji} $ where $E_{ij}$ denotes the usual
elementary matrix.
The Yangian $Y(\mathfrak{gl}_N)$ is the unital associative algebra over $\mathbb{C}[\hbar]$ generated~by
elements~$t_{ij}^{(r)}$ for~$r\in\mathbb{Z}_+$, $i,j \in \{1,\ldots,N\}$ satisfying $t_{ij}^{(0)}=\delta_{ij}$, together
with the following relations: if $t_{ij}(u)=\sum\limits_{r=0}^\infty {t_{ij}^{(r)}u^{-r}} \in
Y(\mathfrak{gl}_N)[[u^{-1}]]$ and $t(u) =\sum\limits_{i,j=1}^N {E_{ij} \otimes t_{ij}(u)}$, then
\begin{gather}
R(u-v)t_1(u)t_2(v)=t_2(v)t_1(u)R(u-v),
\label{RTT}
\end{gather}
where $t_1(u)$ (resp.\ $t_2(u)$) is the element of $\End_{\mathbb{C}}\big(\mathbb{C}^N\big) \otimes_{\mathbb{C}}
\End_{\mathbb{C}}\big(\mathbb{C}^N\big) \otimes_{\mathbb{C}} Y(\mathfrak{gl}_N)[[u^{-1}]]$ obtained by replacing
$E_{ij}$ by $E_{ij} \otimes I$ (resp.
$I \otimes E_{ij}$) in $t(u)$.
Here, we have also identif\/ied $R(u-v)$ with $R(u-v)\otimes 1$.
Relation~\eqref{RTT} is equivalent to the system of relations
\begin{gather*}
\big[t_{ij}^{(m+1)},t_{kl}^{(n)}\big]-\big[t_{ij}^{(m)},t_{kl}^{(n+1)}\big]
= \hbar\big(t_{kj}^{(m)}t_{il}^{(n)}-t_{kj}^{(n)}t_{il}^{(m)}\big).
\end{gather*}
\end{definition}

\begin{definition}
The quantum af\/f\/ine~$R$-matrix is the element of $\End_{\mathbb{C}}\big(\mathbb{C}^N\big) \otimes
\End_{\mathbb{C}}\big(\mathbb{C}^N\big) \otimes_{\mathbb{C}} \mathbb{C}[u,v]$ given~by
\begin{gather*}
R_q(u,v) = \sum\limits_{i,j=1}^N \big(uq^{-\delta_{ij}} - vq^{\delta_{ij}}\big) E_{ii} \otimes E_{jj} - \big(q-q^{-1}\big)u
\sum\limits_{\substack{i,j=1\\ i > j}}^N E_{ij} \otimes E_{ji}\\
\hphantom{R_q(u,v) =}{}
- \big(q-q^{-1}\big)v \sum\limits_{\substack{i,j=1\\ i < j}}^N E_{ij} \otimes E_{ji}.
\end{gather*}
\end{definition}
Set $\mathcal{L}(\mathfrak{gl}_N)=\mathfrak{gl}_N \otimes_{\mathbb{C}} \mathbb{C}[s,s^{-1}]$.
\begin{definition}
The quantum loop algebra $\mathfrak{U}_q(\mathcal{L}(\mathfrak{gl}_N))$ is the unital associative algebra
over~$\mathbb{C}(q)$ generated by $\{T_{ij}^{(r)},\overline{T}_{ij}^{(r)} \,|\, r\in\mathbb{Z}_+, \; i,j=1,\ldots,N\}$
which must satisfy the following relations
\begin{gather}
T_{ij}^{(0)} = 0 = \overline{T}_{ji}^{(0)} \qquad \text{if} \qquad  1 \leq i<j \leq N,
\nonumber
\\
T_{ii}^{(0)} \overline{T}_{ii}^{(0)} = 1 = \overline{T}_{ii}^{(0)} T_{ii}^{(0)} \qquad \forall\,   1 \leq i \leq N,
\nonumber
\\
R_q(u,v)T_1(u)T_2(v) = T_2(v) T_1(u) R_q(u,v),
\label{RTTql1}
\\
R_q(u,v)\overline{T}_1(u)\overline{T}_2(v) = \overline{T}_2(v) \overline{T}_1(u) R_q(u,v),
\label{RTTql2}
\\
R_q(u,v)\overline{T}_1(u) T_2(v) = T_2(v) \overline{T}_1(u) R_q(u,v),
\label{RTTql3}
\end{gather}
$T_a(u)$ and $\overline{T}_a(u)$ are obtained from $T(u):=\sum\limits_{i,j=1}^N{E_{ij} \otimes T_{ij}(u)}$ and
$\overline{T}(u):=\sum\limits_{i,j=1}^N{E_{ij} \otimes \overline{T}_{ij}(u)}$ in the same way as $t_a(u)$ in
Def\/inition~\ref{Ygln}; in this case, $T_{ij}(u)$, $\overline{T}_{ij}(u)$ are the elements
\begin{gather*}
T_{ij}(u)=\sum\limits_{r=0}^\infty{T_{ij}^{(r)}u^{-r}} \in
\mathfrak{U}_q(\mathcal{L}(\mathfrak{gl}_N))\big[\big[u^{-1}\big]\big],
\qquad
\overline{T}_{ij}(u)=\sum\limits_{r=0}^\infty{\overline{T}_{ij}^{(r)}u^{r}} \in
\mathfrak{U}_q(\mathcal{L}(\mathfrak{gl}_N))[[u]].
\end{gather*}
\end{definition}

Relations~\eqref{RTTql1},\eqref{RTTql2} and~\eqref{RTTql3} can be made more explicit in terms of the generators of
$\mathfrak{U}_q(\mathcal{L}(\mathfrak{gl}_N))$.
For instance,~\eqref{RTTql1} is equivalent to the following family of relations
\begin{gather}
\big(q^{-\delta_{ik}} T_{ij}^{(r+1)} T_{kl}^{(s)} - q^{\delta_{ik}}T_{ij}^{(r)} T_{kl}^{(s+1)}\big) - \big(q^{-\delta_{jl}}
T_{kl}^{(s)} T_{ij}^{(r+1)} - q^{\delta_{jl}} T_{kl}^{(s+1)} T_{ij}^{(r)}\big)
\nonumber
\\
\qquad{}
= \big(q-q^{-1}\big) \big(\delta_{i>k} T_{kj}^{(r+1)} T_{il}^{(s)} + \delta_{i<k} T_{kj}^{(r)} T_{il}^{(s+1)}\big)
\nonumber\\
\qquad\phantom{=}
- \big(q-q^{-1}\big)\big(\delta_{l>j} T_{kj}^{(s)} T_{il}^{(r+1)} + \delta_{l<j} T_{kj}^{(s+1)} T_{il}^{(r)}\big).
\label{RTTql12}
\end{gather}

Let $\mathcal{A}$ be the localization of $\mathbb{C}[q,q^{-1}]$ at the ideal $(q-1)$.
Let $\mathfrak{U}_{\mathcal{A}}(\mathcal{L}(\mathfrak{gl}_N))$ be the $\mathcal{A}$-subalgebra of
$\mathfrak{U}_q(\mathcal{L}(\mathfrak{gl}_N))$ generated by the elements $ \tau_{ij}^{(r)}$,
$\overline{\tau}_{ij}^{(r)}$ given~by
\begin{gather*}
\tau_{ij}^{(r)} = \frac{T_{ij}^{(r)}}{q-q^{-1}},
\qquad
\overline{\tau}_{ij}^{(r)} = \frac{\overline{T}_{ij}^{(r)}}{q-q^{-1}}
\qquad\text{for}\quad r\ge 0,\quad 1\le i,j \le N,
\end{gather*}
except that, when $r=0$ and $i=j$, we set
\begin{gather*}
\tau_{ii}^{(0)} = \frac{T_{ii}^{(0)}-1}{q-1}, \qquad \overline{\tau}_{ii}^{(0)} = \frac{\overline{T}_{ii}^{(0)}-1}{q-1}.
\end{gather*}

\begin{theorem}[Section~3 of~\cite{MRS}]
\label{qlim}
The assignment $E_{ij} s^r \mapsto \tau_{ij}^{(r)}$ $\forall\, r\ge 0$, $1\le i$, $j\le n$ except if $r=0$ and $1\le i<j\le n$,
$E_{ij} s^{-r} \mapsto -\overline{\tau}_{ij}^{(r)}$ $\forall\, r\ge 0$, $1\le i$, $j\le n$ except if $r=0$ and $1\le j<i\le n$,
induces an isomorphism $\mathfrak{U}(\mathcal{L}(\mathfrak{gl}_N)) \stackrel{\sim}{\longrightarrow}
\mathfrak{U}_{\mathcal{A}}(\mathcal{L}(\mathfrak{gl}_N)) \otimes_{\mathcal{A}} \mathbb{C}$, where $\mathbb{C}$
is viewed as an $\mathcal{A}$-module via $\mathcal{A}/(q-1) \stackrel{\sim}{\longrightarrow} \mathbb{C}$.
\end{theorem}

We have the following sequence of algebra homomorphisms similar to the one considered in~\cite{GuMa}
\begin{gather*}
\mathfrak{U}_{\mathcal{A}}(\mathcal{L}(\mathfrak{gl}_N)) \twoheadrightarrow
\mathfrak{U}_{\mathcal{A}}(\mathcal{L}(\mathfrak{gl}_N)) /
(q-1)\mathfrak{U}_{\mathcal{A}}(\mathcal{L}(\mathfrak{gl}_N)) \stackrel{\sim}{\longrightarrow}
\mathfrak{U}(\mathcal{L}(\mathfrak{gl}_N)) \stackrel{s\mapsto 1}{\twoheadrightarrow}
\mathfrak{U}(\mathfrak{gl}_N).
%\label{evalhom}
\end{gather*}



For $m\ge 0$, denote by $\mathsf{K}_m$ the Lie ideal of $\mathcal{L}(\mathfrak{gl}_N)$
spanned by $X\otimes s^r (s-1)^m \forall\, r\in\mathbb{Z}$, $\forall\, X\in\mathfrak{gl}_N$.
Let~$U$ be the subspace of $\mathfrak{U}_{\mathcal{A}}(\mathcal{L}(\mathfrak{gl}_N))$ spanned over
$\mathbb{C}$ by the generators $\tau_{ij}^{(r)}$, $\overline{\tau}_{ij}^{(r)}$ and observe that
 \begin{gather*}
 U \cap (q-1){\mathfrak{U}}_{\mathcal{A}}({\mathcal{L}}({\mathfrak{gl}}_N))
 = \operatorname{span}_{\mathbb{C}} \big\{\tau_{ii}^{(0)}
 +\overline{\tau}_{ii}^{(0)} \,\big|\, i=1,\ldots,N\big\}.
\end{gather*}
This is because $(q-1)\mathfrak{U}_{\mathcal{A}}(\mathcal{L}(\mathfrak{gl}_N))=\Ker(\psi)$,
where~$\psi$ is the composite
\begin{gather*}
\mathfrak{U}_{\mathcal{A}}(\mathcal{L}(\mathfrak{gl}_N)) \twoheadrightarrow
\mathfrak{U}_{\mathcal{A}}(\mathcal{L}(\mathfrak{gl}_N)) /
(q-1)\mathfrak{U}_{\mathcal{A}}(\mathcal{L}(\mathfrak{gl}_N)) \stackrel{\sim}{\longrightarrow}
\mathfrak{U}(\mathcal{L}(\mathfrak{gl}_N)).
\end{gather*}

Moreover, notice that
\begin{gather*}
\tau_{ii}^{(0)}+\overline{\tau}_{ii}^{(0)} = \frac{\overline{T}_{ii}^{(0)}\big(T_{ii}^{(0)}-1\big)^2}{q-1} = (q-1)
\big((q-1)\overline{\tau}_{ii}^{(0)}+1\big) \big(\tau_{ii}^{(0)}\big)^2 \in
(q-1)\mathfrak{U}_{\mathcal{A}}(\mathcal{L}(\mathfrak{gl}_N)).
\end{gather*}
Let $\mathbf{K}_0 = \mathfrak{U}_{\mathcal{A}}(\mathcal{L}(\mathfrak{gl}_N))$.
For $m \geq 1$, let $\mathbb{K}_m$ be the two-sided ideal of
$\mathfrak{U}_{\mathcal{A}}(\mathcal{L}(\mathfrak{gl}_N))$ generated by $\psi^{-1}(\mathsf{K}_m) \cap U$, and
set $\mathbf{K}_m$ to be the sum of the ideals $(q-q^{-1})^{m_0}\mathbb{K}_{m_1} \cdots \mathbb{K}_{m_k}$ with $m_0 + m_1
+ \cdots + m_k \ge m$.
This is slightly dif\/ferent from the def\/inition of the analogous ideals $\mathbf{K}_m$ in~\cite{GuMa} in the case of
$\mathfrak{s}\mathfrak{l}_N$ because, for the $\mathfrak{gl}_N$ case, $\mathbf{K}_1^m$ is strictly smaller
than $\mathbf{K}_m$.

Let $\widetilde{Y}(\mathfrak{gl}_N)$ be the $\mathbb{C}$-algebra
\begin{gather*}
\widetilde{Y}(\mathfrak{gl}_N) = \bigoplus_{m=0}^{\infty} \mathbf{K}_{m}/\mathbf{K}_{m+1}.
\end{gather*}
$\widetilde{Y}(\mathfrak{gl}_N)$ can be viewed as a~$\mathbb{C}[\hbar]$-algebra if we set
$\hbar=\overline{q-q^{-1}} \in \mathbf{K}_1/\mathbf{K}_2$.
Note that the f\/irst quotient is isomorphic to $\mathfrak{U}(\mathfrak{gl}_N)$ by def\/inition.

\begin{theorem}
\label{mainthm}
$ \widetilde{Y}(\mathfrak{gl}_N)$ is isomorphic to $Y(\mathfrak{gl}_N)$.
\end{theorem}

For the analogue of our f\/irst theorem for an arbitrary complex semisimple Lie algebra, see~\cite{Dr1} and~\cite{GuMa}.

For $m,r\ge 0$, def\/ine recursively elements $T_{ij}^{(r,m)}$ in the following way
\begin{gather*}
T_{ij}^{(r,0)} = \tau_{ij}^{(r)} \qquad\text{and}\qquad T_{ij}^{(r,m+1)} = T_{ij}^{(r+1,m)} - T_{ij}^{(r,m)},
\end{gather*}
except that if $i<j$, $T_{ij}^{(0,0)} = -\overline{\tau}_{ij}^{(0)}$.

{\sloppy One can check by induction on~$m$ that, for every~$r$, $\psi\big(T_{ij}^{(r,m)}\big)=E_{ij}s^r(s-1)^m$, hence $T_{ij}^{(r,m)}
\in \mathbf{K}_m$.
Set $\xi_{ij}^{(r,m)} = \overline{T_{ij}^{(r,m)}} \in \mathbf{K}_m/\mathbf{K}_{m+1}$.

}


\begin{proof}[Proof of Theorem~\ref{mainthm}] We will prove that there is an isomorphism $\varphi\colon Y(\mathfrak{gl}_N)
\stackrel{\sim}{\longrightarrow} \widetilde{Y}(\mathfrak{gl}_N)$ given by $t_{ij}^{(m+1)} \mapsto
\xi_{ij}^{(0,m)}$ for $m\ge 0$.

Relation~\eqref{RTTql12} can be rewritten in the following way
\begin{gather*}
q^{-\delta_{ik}}\big(T_{ij}^{(r+1)} - T_{ij}^{(r)}\big) T_{kl}^{(s)}- T_{ij}^{(r)} \big(q^{\delta_{ik}}T_{kl}^{(s+1)}
- q^{-\delta_{ik}}T_{kl}^{(s)}\big)
\\
\qquad\phantom{=}{}
- \big(q^{-\delta_{jl}}T_{kl}^{(s)} \big(T_{ij}^{(r+1)} - T_{ij}^{(r)}\big) - \big(q^{\delta_{jl}}T_{kl}^{(s+1)} -
q^{-\delta_{jl}}T_{kl}^{(s)}\big) T_{ij}^{(r)} \big)
\\
\qquad
= \big(q-q^{-1}\big) \big(\delta_{i>k}T_{kj}^{(r+1)} T_{il}^{(s)} + \delta_{i<k} T_{kj}^{(r)} T_{il}^{(s+1)}\big)
\\
\qquad\phantom{=}{}
  - \big(q-q^{-1}\big) \big(\delta_{l>j} T_{kj}^{(s)} T_{il}^{(r+1)} + \delta_{l<j} T_{kj}^{(s+1)} T_{il}^{(r)} \big).
\end{gather*}

If $r,s \geq 1$ then, after rearranging and dividing both sides by $(q-q^{-1})^2$, we get
\begin{gather*}
q^{-\delta_{ik}} \big(T_{ij}^{(r,1)} T_{kl}^{(s,0)}- T_{ij}^{(r,0)} T_{kl}^{(s,1)} \big) - \big(q^{\delta_{ik}} -
q^{-\delta_{ik}}\big) T_{ij}^{(r,0)} T_{kl}^{(s+1,0)}
\\
\qquad\phantom{=}{}
 - q^{-\delta_{jl}} \big(T_{kl}^{(s,0)} T_{ij}^{(r,1)} - T_{kl}^{(s,1)} T_{ij}^{(r,0)} \big) + \big(q^{\delta_{jl}} -
q^{-\delta_{jl}}\big) T_{kl}^{(s+1,0)} T_{ij}^{(r,0)}
\\
\qquad
= \big(q-q^{-1}\big) \big(  \delta_{i>k} T_{kj}^{(r+1,0)} T_{il}^{(s,0)} + \delta_{i<k} T_{kj}^{(r,0)} T_{il}^{(s+1,0)}\big)
\\
\qquad\phantom{=}{}
  - \big(q-q^{-1}\big) \big(\delta_{l>j} T_{kj}^{(s,0)} T_{il}^{(r+1,0)} + \delta_{l<j} T_{kj}^{(s+1,0)} T_{il}^{(r,0)} \big).
\end{gather*}

Using $T_{ij}^{(r+1,m)} - T_{ij}^{(r,m)}=T_{ij}^{(r,m+1)}$ and $T_{kl}^{(s+1,n)} - T_{kl}^{(s,n)}=T_{kl}^{(s,n+1)}$, we
deduce by induction on~$m$ and~$n$ that, for all $r,s \geq 1$ and all $m,n\ge 0$,
\begin{gather}
q^{-\delta_{ik}} \big(T_{ij}^{(r,m+1)}T_{kl}^{(s,n)} - T_{ij}^{(r,m)} T_{kl}^{(s,n+1)} \big) - \big(q^{\delta_{ik}} -
q^{-\delta_{ik}}\big) T_{ij}^{(r,m)} T_{kl}^{(s+1,n)}
\nonumber
\\
\qquad\phantom{=}{}
  - q^{-\delta_{jl}} \big(T_{kl}^{(s,n)} T_{ij}^{(r,m+1)} - T_{kl}^{(s,n+1)} T_{ij}^{(r,m)} \big) + \big(q^{\delta_{jl}} -
q^{-\delta_{jl}}\big) T_{kl}^{(s+1,n)} T_{ij}^{(r,m)}
\nonumber
\\
\qquad
= \big(q-q^{-1}\big) \big(  \delta_{i>k} T_{kj}^{(r+1,m)} T_{il}^{(s,n)} + \delta_{i<k} T_{kj}^{(r,m)} T_{il}^{(s+1,n)}\big)
\nonumber
\\
\qquad\phantom{=}{}
  - \big(q-q^{-1}\big) \big(\delta_{l>j} T_{kj}^{(s,n)} T_{il}^{(r+1,m)} + \delta_{l<j} T_{kj}^{(s+1,n)} T_{il}^{(r,m)} \big).
\label{rs}
\end{gather}
Consider the case $r=s=1$ in~\eqref{rs}.
Using $T_{ij}^{(r+1,m)}=T_{ij}^{(r,m+1)}+T_{ij}^{(r,m)}$ and $T_{kl}^{(r+1,n)}=T_{kl}^{(r,n+1)}+T_{kl}^{(r,n)}$ for all
$r\ge 0$, we obtain, for all $m,n\ge 0$,
\begin{gather*}
  q^{-\delta_{ik}} \big(T_{ij}^{(0,m+2)} + T_{ij}^{(0,m+1)}\big) T_{kl}^{(0,n+1)} -
q^{-\delta_{ik}}\big(T_{ij}^{(0,m+1)} + T_{ij}^{(0,m)}\big) T_{kl}^{(0,n+2)}
\\
\qquad\phantom{=}{}
  - \big(q^{\delta_{ik}} - q^{-\delta_{ik}}\big) \big(T_{ij}^{(0,m+1)} + T_{ij}^{(0,m)}\big) T_{kl}^{(1,n+1)} + q^{-\delta_{ik}}
\big(T_{ij}^{(0,m+2)} + T_{ij}^{(0,m+1)}\big) T_{kl}^{(0,n)}
\\
\qquad\phantom{=}{}
  - q^{-\delta_{ik}}\big(T_{ij}^{(0,m+1)} + T_{ij}^{(0,m)}\big) T_{kl}^{(0,n+1)} - \big(q^{\delta_{ik}} -
q^{-\delta_{ik}}\big)\big(T_{ij}^{(0,m+1)} + T_{ij}^{(0,m)}\big) T_{kl}^{(1,n)}
\\
\qquad\phantom{=}{}
  - q^{-\delta_{jl}} T_{kl}^{(0,n+1)} \big(T_{ij}^{(0,m+2)} + T_{ij}^{(0,m+1)}\big) +
q^{-\delta_{jl}}T_{kl}^{(0,n+2)}\big(T_{ij}^{(0,m+1)} + T_{ij}^{(0,m)}\big)
\\
\qquad\phantom{=}{}
  + \big(q^{\delta_{jl}} - q^{-\delta_{jl}}\big) T_{kl}^{(1,n+1)} \big(T_{ij}^{(0,m+1)} + T_{ij}^{(0,m)}\big) - q^{-\delta_{jl}}
T_{kl}^{(0,n)} \big(T_{ij}^{(0,m+2)} + T_{ij}^{(0,m+1)}\big)
\\
\qquad\phantom{=}{}
  + q^{-\delta_{jl}}T_{kl}^{(0,n+1)} \big(T_{ij}^{(0,m+1)} + T_{ij}^{(0,m)}\big) + \big(q^{\delta_{jl}} - q^{-\delta_{jl}}\big)
T_{kl}^{(1,n)} \big(T_{ij}^{(0,m+1)} + T_{ij}^{(0,m)}\big)
\\
\qquad
  = \big(q-q^{-1}\big) \Big(\delta_{i>k} \big(T_{kj}^{(1,m+1)} + T_{kj}^{(1,m)}\big) T_{il}^{(0,n+1)} + \delta_{i<k}
\big(T_{kj}^{(0,m+1)} + T_{kj}^{(0,m)}\big) T_{il}^{(1,n+1)} \Big)
\\
\qquad\phantom{=}{}
+ \big(q-q^{-1}\big) \Big(\delta_{i>k} \big(T_{kj}^{(1,m+1)} + T_{kj}^{(1,m)}\big) T_{il}^{(0,n)} + \delta_{i<k}
\big(T_{kj}^{(0,m+1)} + T_{kj}^{(0,m)}\big) T_{il}^{(1,n)}\Big)
\\
\qquad\phantom{=}{}
- \big(q-q^{-1}\big) \Big(\delta_{l>j} T_{kj}^{(0,n+1)} \big(T_{il}^{(1,m+1)} + T_{il}^{(1,m)}\big) + \delta_{l<j}
T_{kj}^{(1,n+1)} \big(T_{il}^{(0,m+1)} + T_{il}^{(0,m)}\big) \Big)
\\
\qquad\phantom{=}{}
- \big(q-q^{-1}\big) \Big(\delta_{l>j} T_{kj}^{(0,n)} \big(T_{il}^{(1,m+1)} + T_{il}^{(1,m)}\big) + \delta_{l<j} T_{kj}^{(1,n)}
\big(T_{il}^{(0,m+1)} + T_{il}^{(0,m)}\big) \Big).
\end{gather*}

Notice that both sides of this last equality are in $\mathbf{K}_{m+n+1}$ (and some of the terms are even in
$\mathbf{K}_{m+n+2}$ or in $\mathbf{K}_{m+n+3}$).
Modulo $\mathbf{K}_{m+n+2}$, we obtain the congruence
\begin{gather*}
q^{-\delta_{ik}} \big(T_{ij}^{(0,m+1)} T_{kl}^{(0,n)}- T_{ij}^{(0,m)} T_{kl}^{(0,n+1)} \big) - \big(q^{\delta_{ik}} -
q^{-\delta_{ik}}\big) T_{ij}^{(0,m)} T_{kl}^{(1,n)}
\\
\qquad\phantom{=}{}
  - q^{-\delta_{jl}} \big(T_{kl}^{(0,n)} T_{ij}^{(0,m+1)} - T_{kl}^{(0,n+1)} T_{ij}^{(0,m)} \big) + \big(q^{\delta_{jl}} -
q^{-\delta_{jl}}\big) T_{kl}^{(1,n)} T_{ij}^{(0,m)}
\\
\qquad
\equiv \big(q-q^{-1}\big) \big(  \delta_{i>k} T_{kj}^{(1,m)} T_{il}^{(0,n)} + \delta_{i<k} T_{kj}^{(0,m)} T_{il}^{(1,n)}\big)
\\
\qquad\phantom{=}{}
  - \big(q-q^{-1}\big) \big(\delta_{l>j} T_{kj}^{(0,n)} T_{il}^{(1,m)} + \delta_{l<j} T_{kj}^{(1,n)} T_{il}^{(0,m)}\big).
\end{gather*}
Moreover, modulo $\mathbf{K}_{m+n+2}$, we also have
\begin{gather*}
\big(q^{\delta_{ik}} - q^{-\delta_{ik}}\big) T_{ij}^{(0,m)} T_{kl}^{(1,n)} \equiv \big(q^{\delta_{ik}} - q^{-\delta_{ik}}\big)
T_{ij}^{(0,m)} T_{kl}^{(0,n)},
\\
\big(q^{\delta_{jl}}-q^{-\delta_{jl}}\big) T_{kl}^{(1,n)} T_{ij}^{(0,m)} \equiv \big(q^{\delta_{jl}}-q^{-\delta_{jl}}\big)
T_{kl}^{(0,n)} T_{ij}^{(0,m)},
\\
\big(q-q^{-1}\big) T_{kj}^{(1,m)} T_{il}^{(0,n)} \equiv \big(q-q^{-1}\big)T_{kj}^{(0,m)} T_{il}^{(0,n)},
\\
\big(q-q^{-1}\big) T_{kj}^{(0,m)}
T_{il}^{(1,n)} \equiv \big(q-q^{-1}\big) T_{kj}^{(0,m)} T_{il}^{(0,n)},
\\
\big(q-q^{-1}\big) T_{kj}^{(0,n)} T_{il}^{(1,m)} \equiv \big(q-q^{-1}\big) T_{kj}^{(0,n)} T_{il}^{(0,m)},
\\
\big(q-q^{-1}\big) T_{kj}^{(1,n)}
T_{il}^{(0,m)} \equiv \big(q-q^{-1}\big) T_{kj}^{(0,n)} T_{il}^{(0,m)}.
\end{gather*}
Therefore, passing to the quotient $\mathbf{K}_{m+n+1}/\mathbf{K}_{m+n+2}$, we obtain
\begin{gather*}
\big(\xi_{ij}^{(0,m+1)} \xi_{kl}^{(0,n)} - \xi_{ij}^{(0,m)} \xi_{kl}^{(0,n+1)} \big)- \delta_{ik} \hbar
\xi_{ij}^{(0,m)} \xi_{kl}^{(0,n)}
\\
\qquad\phantom{=}{}
  - \big(\xi_{kl}^{(0,n)} \xi_{ij}^{(0,m+1)} - \xi_{kl}^{(0,n+1)} \xi_{ij}^{(0,m)} \big) + \delta_{jl} \hbar
\xi_{kl}^{(0,n)} \xi_{ij}^{(0,m)}
\\
\qquad
= \hbar \big(\delta_{i>k} \xi_{kj}^{(0,m)} \xi_{il}^{(0,n)}+ \delta_{i<k} \xi_{kj}^{(0,m)} \xi_{il}^{(0,n)}\big)
  - \hbar \big(\delta_{l>j} \xi_{kj}^{(0,n)} \xi_{il}^{(0,m)} + \delta_{l<j} \xi_{kj}^{(0,n)} \xi_{il}^{(0,m)}\big).
\end{gather*}
This last relation is equivalent to
\begin{gather*}
\big[\xi_{ij}^{(0,m+1)}, \xi_{kl}^{(0,n)}\big] - \big[\xi_{ij}^{(0,m)}, \xi_{kl}^{(0,n+1)}\big] = \hbar \big(\xi_{kj}^{(0,m)}
\xi_{il}^{(0,n)} - \xi_{kj}^{(0,n)} \xi_{il}^{(0,m)}\big).
\end{gather*}
This holds for all $m,n\ge 0$.

All the previous computations prove that $\varphi\colon Y(\mathfrak{gl}_N) \longrightarrow
\widetilde{Y}(\mathfrak{gl}_N)$ given by $\varphi(t_{ij}^{(m+1)}) = \xi_{ij}^{(0,m)}$ for $m\ge 0$ is an
algebra homomorphism.
We still have to show that~$\varphi$ is injective and surjective.

We will f\/irst demonstrate surjectivity.
Towards this end, we def\/ine elements $\overline{T}_{ij}^{(r,m)}$ as follows.
Let $\overline{T}_{ij}^{(r,0)}=\overline{\tau}_{ij}^{(r)}$, except that $\overline{T}_{ij}^{(0,0)}=-\tau_{ij}^{(0)}$
when $i \geq j$.
Then, for each $m \geq 0$, let
\begin{gather*}
\overline{T}_{ij}^{(r,m+1)}=\overline{T}_{ij}^{(r+1,m)}-\overline{T}_{ij}^{(r,m)}.
\end{gather*}

Also for each $m \geq 0$, let $\widetilde{T}_{ij}^{(0,m)}=\overline{T}_{ij}^{(0,m)}$ and
$\widetilde{T}_{ij}^{(m,m)}=(-1)^{m+1}T_{ij}^{(0,m)}$, and for $1 \leq r \leq m$ def\/ine recursively
\begin{gather*}
\widetilde{T}_{ij}^{(r,m+1)}=\widetilde{T}_{ij}^{(r-1,m)}-\widetilde{T}_{ij}^{(r,m)}.
\end{gather*}


Induction on~$m$ shows that the elements $T_{ij}^{(r,m)}$, $\overline{T}_{ij}^{(r,m)}$ and $\widetilde{T}_{ij}^{(r,m)}$
respectively map via~$\psi$ to the elements $E_{ij}s^r(s-1)^m$, $(-1)^{m+1}E_{ij}s^{-(m+r)}(s-1)^m$ and
$(-1)^{m+1}E_{ij}s^{-(m-r)}(s-1)^m$ in $\mathfrak{U}(\mathcal{L}(\mathfrak{gl}_N))$.
It follows that, for f\/ixed~$m$, the images of those elements under~$\psi$ span $\mathsf{K}_m$.
Moreover, all those elements are in~$U$ by def\/inition.

Note that for any f\/ixed $X \in \psi^{-1}(\mathsf{K}_m) \cap U$, there exists some element~$Y$ in
 \begin{gather*}
\operatorname{span}_{\mathbb{C}}\big\{T_{ij}^{(r,m)},\overline{T}_{ij}^{(r,m)},\widetilde{T}_{ij}^{(r,m)} \,|\, i,j=1,\ldots,N,\; r \in
\mathbb{Z}_+\big\}
\end{gather*}
for which $X-Y \in (q-1)\mathfrak{U}_{\mathcal{A}}(\mathcal{L}(\mathfrak{gl}_N))$, because the map
$\mathfrak{U}_{\mathcal{A}}\mathcal{L}(\mathfrak{gl}_N)
/(q-1)\mathfrak{U}_{\mathcal{A}}(\mathcal{L}(\mathfrak{gl}_N))
\stackrel{\sim}{\longrightarrow} \mathfrak{U}(\mathcal{L}(\mathfrak{gl}_N))$ is an isomorphism.
Since $X-Y$ is also in~$U$, we have{\samepage
\begin{gather*}
X-Y \in \operatorname{span}_{\mathbb{C}} \big\{\tau_{ii}^{(0)}+\overline{\tau}_{ii}^{(0)} \,|\, i=1,\ldots,N\big\} \subset \mathbf{K}_{\ell}
 \; \forall\, \ell \geq 0.
\end{gather*}
That $\tau_{ii}^{(0)}+\overline{\tau}_{ii}^{(0)}$ is in $\mathbf{K}_{\ell}$ for all $\ell \geq 0$ is a~consequence of the
fact that $\psi(\tau_{ii}^{(0)}+\overline{\tau}_{ii}^{(0)})=0$.}

It follows that any element of $\mathbf{K}_m$ is congruent modulo $\mathbf{K}_{m+1}$ to a~sum of monomials of the form
$f(q)(q-q^{-1})^{m_0}\mathcal{M}$ where $f(q) \in \mathcal{A}$ is not divisible by $q-1$ and $\mathcal{M}=\tau_{i_1
j_1}^{(r_1,m_1)} \cdots \tau_{i_k j_k}^{(r_k,m_k)}$ with
\begin{gather*}
\tau_{i_d j_d}^{(r_d,m_d)} \in \big\{T_{i_d j_d}^{(r_d,m_d)}, \overline{T}_{i_d j_d}^{(r_d,m_d)}, \widetilde{T}_{i_d
j_d}^{(r_d,m_d)}\big\}
\end{gather*}
and $m_0+\cdots+m_k \geq m$.
Moreover, since $T_{i_d j_d}^{(r_d,m_d)}-T_{i_d j_d}^{(r_d-1,m_d)}=T_{i_d j_d}^{(r_d-1,m_d+1)} \in \mathbb{K}_{m_d+1}$
(and similarly for $\overline{T}_{i_d j_d}^{(r_d,m_d)}$ and $\widetilde{T}_{i_d j_d}^{(r_d,m_d)}$), we can reduce modulo
$\mathbf{K}_{m+1}$ to the case when $r_d=0$ for each $d=1,\ldots,k$.

Observe that modulo $\mathbb{K}_{m_d+1}$, we have
\begin{gather}
\overline{T}_{ij}^{(0,m_d)} \equiv (-1)^{m_d+1}T_{ij}^{(0,m_d)} \equiv \widetilde{T}_{ij}^{(0,m_d)}.
\label{TbarT}
\end{gather}
To see this, just take the dif\/ference of the elements on each side (this dif\/ference is in~$U$ by def\/inition) and
apply~$\psi$.
Finally, observe that we can replace $f(q)$ by $f(1)$ modulo $\mathbf{K}_{1}$.

In summary, we have shown that each of the monomials $f(q)(q-q^{-1})^{m_0}\mathcal{M}$ in $\mathbf{K}_m$ is congruent
modulo $\mathbf{K}_{m+1}$ to
\begin{gather*}
f(1)\big(q-q^{-1}\big)^{m_0} T_{i_1 j_1}^{(0,m_1)} \cdots T_{i_k,j_k}^{(0,m_k)}
\end{gather*}
up to a~sign.
The image modulo $\mathbf{K}_{m+1}$ of this element is
\begin{gather*}
f(1)\hbar^{m_0}\xi_{i_1 j_1}^{(0,m_1)} \cdots \xi_{i_k j_k}^{(0,m_k)}
\end{gather*}
and this is in the image of~$\varphi$ by def\/inition.
This completes the proof that~$\varphi$ is surjective.

To prove that~$\varphi$ is injective, it is enough to show that the basis of $Y(\mathfrak{gl}_N)$ given~by
ordered monomials (for some f\/ixed order) in the generators $t_{ij}^{(m)}$ is mapped via~$\varphi$ to some linearly
independent set in $\widetilde{Y}(\mathfrak{gl}_N)$.

By def\/inition, any two ordered monomials $\hbar^{m_0}t_{i_1 j_1}^{(m_1+1)} \cdots t_{i_a j_a}^{(m_a+1)}$ and
$\hbar^{n_0}t_{k_1 l_1}^{(n_1+1)} \cdots t_{k_b l_b}^{(n_b+1)}$ with each $m_d, n_d \geq 0$ and $m_0+m_1+\cdots+m_a \neq
n_0+n_1 + \cdots + n_b$ are mapped via~$\varphi$ to distinct graded pieces in $\widetilde{Y}(\mathfrak{gl}_N)$.
It therefore suf\/f\/ices to show that for each f\/ixed~$m$, the images under~$\varphi$ of all the ordered monomials
$\hbar^{m_0} t_{i_1 j_1}^{(m_1+1)} \cdots t_{i_a j_a}^{(m_a+1)}$ with $m_1+\cdots+m_a=m$ are linearly independent in
$\mathbf{K}_m/\mathbf{K}_{m+1}$.

Consider any linear combination over $\mathbb{C}$ of ordered monomials $(\overline{q-q^{-1}})^{m_0}\xi_{i_1
j_1}^{(0,m_1)} \cdots \xi_{i_a j_a}^{(0,m_a)}$ with $m_0+m_1+\cdots+m_a=m$, and suppose that this sum is zero in
$\mathbf{K}_m/\mathbf{K}_{m+1}$.
Then we have a~linear combination~$S$ of ordered monomials $(q-q^{-1})^{m_0}T_{i_1 j_1}^{(0,m_1)} \cdots T_{i_a
j_a}^{(0,m_a)}$ which is not just in~$\mathbf{K}_m$, but also in~$\mathbf{K}_{m+1}$.
We can assume that the minimum of the exponents~$m_0$ is~$0$.

$\psi(S)$ is a~linear combination of some of the ordered monomials $E_{i_1 j_1}(s-1)^{m_1} \cdots E_{i_a j_a}
(s-1)^{m_a}$.
On the other hand, since~$S$ is also in $\mathbf{K}_{m+1}$, $\psi(S)$ can be expressed as a~linear combination of
monomials of the form $E_{k_1 l_1} s^{r_1} (s-1)^{n_1} \cdots E_{k_b l_b} s^{r_b} (s-1)^{n_b}$ with $r_1,\ldots,r_b \in
\mathbb{Z}$ and $n_1+\cdots+n_b \geq m+1$.
If $r_1 = \cdots = r_b=0$, then this is impossible unless the coef\/f\/icients of both linear combinations vanish.
Let's prove that the same is true more generally.

For each $r \geq 1$, we have a~composite of algebra homomorphisms
\begin{gather*}
\mathfrak{U}(\mathcal{L}(\mathfrak{gl}_N)) \stackrel{\Delta}{\to}
\mathfrak{U}(\mathcal{L}(\mathfrak{gl}_N))^{\otimes r} \stackrel{f^{\otimes r}}{\rightarrow}
\End_{\mathbb{C}}\big(\mathbb{C}^N\big)^{\otimes r} \otimes \mathbb{C}\big[x_1^{\pm 1},x_2^{\pm 1},\ldots,x_r^{\pm 1}\big],
\end{gather*}
where~$\Delta$ is the standard coproduct on the enveloping algebra of a~Lie algebra and the homomorphism $f\colon
\mathfrak{U}(\mathcal{L}(\mathfrak{gl}_N)) \to \End_{\mathbb{C}}\big(\mathbb{C}^N\big) \otimes
\mathbb{C}[x,x^{-1}]$ is given by $f(E_{ij}x^t) = E_{ij} \otimes x^t$.

We also have for each choice of~$r$ nonnegative integers $\alpha_1, \ldots, \alpha_r$ a~dif\/ferential operator
\begin{gather*}
\partial_{\alpha_1,\ldots,\alpha_r}\colon \ \End_{\mathbb{C}}\big(\mathbb{C}^N\big)^{\otimes r} \otimes \mathbb{C}\big[x_1^{\pm
1},x_2^{\pm 1},\ldots,x_r^{\pm 1}\big] \to \End_{\mathbb{C}}\big(\mathbb{C}^N\big)^{\otimes r}
\end{gather*}
given~by
\begin{gather*}
\partial_{\alpha_1, \ldots, \alpha_r}=\left.\frac{\partial^{\alpha_1}}{\partial x_1^{\alpha_1}}
\frac{\partial^{\alpha_2}}{\partial x_2^{\alpha_2}} \cdots \frac{\partial^{\alpha_r}}{\partial x_r^{\alpha_r}}
\right\rvert_{x_1,\ldots,x_r=1}.
\end{gather*}
Take $r \geq \max\{a,b\}$ where the maximum is taken over all the monomials in $\psi(S)$ (with~$a$ and~$b$ related to
the monomials in~$S$ and $\psi(S)$ as above) and note that for any choice of $\alpha_1,\ldots,\alpha_r$ with $\alpha_1+
\cdots + \alpha_r=m$, $\psi(S)$ is in the kernel of the composite $\partial_{\alpha_1,\ldots,\alpha_r} \circ f^{\otimes
r} \circ \Delta$ because $n_1 + \cdots + n_b \geq m+1$.

On the other hand, if~$S$ is nonzero, then we can f\/ind some $\alpha_1,\ldots,\alpha_r$ such that $\alpha_1 + \cdots +
\alpha_r=m$ and $\psi(S)$ is not in the kernel of $\partial_{\alpha_1,\ldots,\alpha_r} \circ f^{\otimes r} \circ
\Delta$: just choose any of the ordered monomials $E_{i_1 j_1}(s-1)^{m_1} \cdots E_{i_a j_a} (s-1)^{m_a}$ in $\psi(S)$
and set $\alpha_1=m_1, \ldots, \alpha_a=m_a$ and $\alpha_d=0$ for $d>a$.

We have just obtained a~contradiction, so $S=0$ and the linear sum of ordered monomials $\xi_{i_1 j_1}^{(0,m_1)} \cdots
\xi_{i_a j_a}^{(0,m_a)}$ must in fact be trivial, as desired.
\end{proof}

\section[Twisted quantum loop algebras and Yangians of type $\mathfrak{o}_N$ and $\mathfrak{sp}_N$]{Twisted quantum loop algebras and Yangians\\ of type $\boldsymbol{\mathfrak{o}_N}$ and $\boldsymbol{\mathfrak{sp}_N}$}
\label{twsec}
We will now prove an analogue of Theorem~\ref{mainthm} for certain twisted Yangians and twisted quantum loop algebras
associated to the symmetric pairs $(\mathfrak{gl}_N,\mathfrak{o}_N)$ and
$(\mathfrak{gl}_N,\mathfrak{sp}_N)$.
We will treat these two cases simultaneously and denote by $\mathfrak{g}_N$ either $\mathfrak{o}_N$ or
$\mathfrak{sp}_N$.
Whenever we use the symbols $\pm$ or $\mp$, it is understood that the sign on the top is used for
$\mathfrak{g}_N=\mathfrak{o}_N$ and the sign on the bottom is used for $\mathfrak{g}_N=\mathfrak{sp}_N$.
In this section, we will use~$t$ to denote transposition in the f\/irst factor of a~tensor product of matrices.

In the orthogonal case, let $G=(g_{ij})$ be the $N \times N$ identity matrix.
For the symplectic case, we take
\begin{gather*}
G=\sum\limits_{k=1}^{N/2} (E_{2k-1,2k}-E_{2k,2k-1}),
\end{gather*}
which makes sense since $\mathfrak{sp}_N$ is only def\/ined when~$N$ is even.
Similarly, let $B=(b_{ij})$ be the $N \times N$ identity matrix in the orthogonal case, while in the symplectic case we
take
\begin{gather*}
B=\sum\limits_{k=1}^{N/2} (qE_{2k-1,2k}-E_{2k,2k-1}).
\end{gather*}

\begin{definition}
\label{Ytw}
Let $R(u)$ be as given in Def\/inition~\ref{Ygln}.
The twisted Yangian $Y^{tw}(\mathfrak{g}_N)$ is the unital associative algebra over $\mathbb{C}[\hbar]$ generated~by
$\big\{s_{ij}^{(r)} \,|\, r\in\mathbb{Z}_+,\; i,j=1,\ldots,N\big\}$ where $s_{ij}^{(0)}=g_{ij}$, together with the following
relations: if $s_{ij}(u)=\sum\limits_{r=0}^\infty s_{ij}^{(r)}u^{-r}$ and $s(u)=\sum\limits_{i,j=1}^N E_{ij} \otimes
s_{ij}(u)$, then
\begin{gather*}
R(u-v)s_1(u)R^t(-u-v)s_2(v)=s_2(v)R^t(-u-v)s_1(u)R(u-v)
\end{gather*}
and
\begin{gather*}
s^t(-u)= \pm s(u) + \hbar\frac{s(u)-s(-u)}{2u},
\end{gather*}
where $s_1(u)$ (resp.\
$s_2(u)$) is the element of $\End_{\mathbb{C}}\big(\mathbb{C}^N\big) \otimes \End_{\mathbb{C}}\big(\mathbb{C}^N\big) \otimes
Y^{tw}(\mathfrak{g}_N)[[u^{-1}]]$ obtained by replacing $E_{ij}$ by $E_{ij} \otimes I$ (resp.
$I \otimes E_{ij}$) in $s(u)$.
\end{definition}

The twisted Yangian is a~deformation of the universal enveloping algebra of the twisted current algebra
$\mathfrak{g}_N^{tw}[s]$ which is def\/ined in the following way.

\begin{definition}
%\label{twcurrent}
Let~$\sigma$ be the automorphism of $\mathfrak{gl}_N$ given~by
\begin{gather}
\label{orthauto}
\sigma(E_{ij})=-E_{ji}
\end{gather}
if $\mathfrak{g}_N=\mathfrak{o}_N$, while
\begin{gather}
\label{sympauto}
\sigma(E_{ij})=(-1)^{i+j-1}E_{j'i'}
\end{gather}
if $\mathfrak{g}_N=\mathfrak{sp}_N$.
Here, $i'=i-1$ if~$i$ is even and $i'=i+1$ is~$i$ is odd.
The twisted current algebra is the subalgebra of $\mathfrak{gl}_N[s]$ given~by
\begin{gather*}
\mathfrak{g}_N^{tw}[s] = \{A(s) \in \mathfrak{gl}_N[s] \,|\, \sigma(A(s))=A(-s)\}.
\end{gather*}
\end{definition}

The twisted Yangians can be regarded as subalgebras of the Yangian for $\mathfrak{gl}_N$:

\begin{proposition}[\cite{Ol}]
\label{embY}
The assignment
\begin{gather*}
s_{ij}^{(r)} \mapsto \sum\limits_{k=1}^N \left(g_{kj}t_{ik}^{(r)} + (-1)^r g_{ik} t_{jk}^{(r)} \right) + \hbar
\sum\limits_{k,l=1}^N \sum\limits_{p=1}^{r-1} (-1)^{r-p} g_{kl} t_{ik}^{(p)} t_{jl}^{(r-p)}
\end{gather*}
provides an embedding of $Y^{tw}(\mathfrak{g}_N)$ into $Y(\mathfrak{gl}_N)$.
\end{proposition}

\begin{definition}[\cite{MRS}]
\label{Uqon}
The twisted quantum loop algebra $\mathfrak{U}_q(\mathcal{L}^{tw}(\mathfrak{o}_N))$ is the unital associative algebra
over $\mathbb{C}(q)$ generated by $\{S_{ij}^{(r)} \,|\, r \in \mathbb{Z}_+, \, 1 \leq i, j \leq N\}$ which are subject to the
relations
\begin{gather*}
S_{ij}^{(0)}=0 \qquad \text{if} \quad  1 \leq i < j \leq n,
\qquad
S_{ii}^{(0)}=1 \qquad \forall\,  1 \leq i \leq n,
\\
R_q(u,v)S_1(u)R_q^t\big(u^{-1},v\big)S_2(v)=S_2(v)R_q^t\big(u^{-1},v\big)S_1(u)R_q(u,v),
\end{gather*}
where $S_a(u)$ is obtained from $S(u)$ in the same way as in Def\/inition~\ref{Ytw}, except that in this case we have
$S(u):=\sum\limits_{i,j=1}^N E_{ij} \otimes S_{ij}(u)$ and $S_{ij}(u):=\sum\limits_{r=0}^\infty S_{ij}^{(r)}u^{-r}$.
\end{definition}

\begin{definition}[\cite{MRS}]
%\label{Uqspn}
The twisted quantum loop algebra $\mathfrak{U}_q(\mathcal{L}^{tw}(\mathfrak{sp}_N))$ is the unital associative algebra
over $\mathbb{C}(q)$ with generators $\big\{S_{ij}^{(r)} \,|\, r \in \mathbb{Z}_+, 1 \leq i, j \leq N\big\}$ and
$\big\{(S_{ii'}^{(0)})^{-1} \,|\, i=1,3,\ldots,N{-}1\big\}$, which are subject to the relations
\begin{gather*}
S_{ij}^{(0)}=0 \qquad \text{whenever} \quad i<j \quad \text{and} \quad j \neq i'
\\
S_{i'i'}^{(0)}S_{ii}^{(0)}-q^2S_{i'i}^{(0)}S_{ii'}^{(0)}=q^3, \qquad i=1, 3, \ldots, N-1,
\\
S_{ii'}^{(0)}(S_{ii'}^{(0)})^{-1}=(S_{ii'}^{(0)})^{-1}S_{ii'}^{(0)}=1, \qquad i=1, 3, \ldots, N-1,
\\
R_q(u,v)S_1(u)R_q^t\big(u^{-1},v\big)S_2(v)=S_2(v)R_q^t\big(u^{-1},v\big)S_1(u)R_q(u,v),
\end{gather*}
where $S_a(u)$ is def\/ined here in the same way as in Def\/inition~\ref{Uqon}.
\end{definition}

The twisted quantum loop algebra $\mathfrak{U}_q(\mathcal{L}^{tw}(\mathfrak{g}_N))$ is a~deformation of the universal
enveloping algebra of the twisted loop algebra $\mathfrak{g}_N^{tw}[s,s^{-1}]$ which is def\/ined in the following way:

\begin{definition}
The twisted loop algebra $\mathfrak{g}_N^{tw}[s,s^{-1}]$ is the Lie subalgebra of
$\mathcal{L}(\mathfrak{gl}_N)$ given~by
\begin{gather*}
\mathfrak{g}_N^{tw}\big[s,s^{-1}\big]=\big\{A(s) \in \mathcal{L}(\mathfrak{gl}_N) \,|\, \sigma(A(s))=A\big(s^{-1}\big)\big\},
\end{gather*}
where~$\sigma$ is given by~\eqref{orthauto} for $\mathfrak{g}_N=\mathfrak{o}_N$ and by~\eqref{sympauto} for
$\mathfrak{g}_N=\mathfrak{sp}_N$.
This algebra is also denoted by $\mathcal{L}^{tw}(\mathfrak{g}_N)$.
\end{definition}

The twisted quantum loop algebras can be regarded as subalgebras of
$\mathfrak{U}_q(\mathcal{L}(\mathfrak{gl}_N))$:

\begin{proposition}[\cite{MRS}]\label{embU}
The assignment
\begin{gather*}
S_{ij}^{(r)} \mapsto \sum\limits_{k,l=1}^N \sum\limits_{p=0}^r b_{kl} T_{ik}^{(p)} \overline{T}_{jl}^{(r-p)}
\end{gather*}
provides an embedding of $\mathfrak{U}_q(\mathcal{L}^{tw}(\mathfrak{g}_N))$ into
$\mathfrak{U}_q(\mathcal{L}(\mathfrak{gl}_N))$.
\end{proposition}

For each $r>0$, let $S_{ij}^{(r,0)}=\frac{S_{ij}^{(r)}}{q-q^{-1}}$.
When $\mathfrak{g}_N = \mathfrak{o}_N$, we set $S_{ij}^{(0,0)}=\frac{S_{ij}^{(0)}}{q-q^{-1}}$ when $i>j$ and
$S_{ij}^{(0,0)}=-S_{ji}^{(0,0)}$ when $i\le j$; when $\mathfrak{g}_N = \mathfrak{sp}_N$, we set
$S_{ij}^{(0,0)}=\frac{S_{ij}^{(0)}-b_{ij}}{q-q^{-1}}$ when $i\ge j$ or $j=i'$ and $S_{ij}^{(0,0)} = -S_{ji}^{(0,0)}$
when $i<j$ and $j\neq i'$.

We def\/ine inductively $S_{ij}^{(r,m)}$ by
\begin{gather*}
S_{ij}^{(r,m+1)}=S_{ij}^{(r+1,m)}-S_{ij}^{(r,m)}
\end{gather*}
for $m \geq 0$.
\begin{lemma}
Identify $S_{ij}^{(r)}$ with its image under the embeddings of Proposition~{\rm \ref{embU}}.
For any $r>0$ and $m \geq 0$, we have
\begin{gather*}
S_{ij}^{(r,m)}=\sum\limits_{k=1}^N \left(b_{kj} T_{ik}^{(r,m)} + b_{ik} \overline{T}_{jk}^{(r,m)} \right) +
\big(q-q^{-1}\big) \sum\limits_{k,l=1}^N \sum\limits_{p=1}^{r-1} b_{kl} T_{ik}^{(p,0)} \overline{T}_{jl}^{(r-p,m)}
\\
\phantom{S_{ij}^{(r,m)}=}{}
 + \big(q-q^{-1}\big)\!\left(\sum\limits_{\substack{1 \leq k \leq i\\ 1 \leq l \leq N}}\!\left(\frac{q}{q+1}\right)^{\delta_{ik}}\!
b_{kl} T_{ik}^{(0,0)} \overline{T}_{jl}^{(r,m)} + \sum\limits_{\substack{1 \leq k \leq N \\ j \leq l \leq N}}\!
\left(\frac{q}{q+1}\right)^{\delta_{jl}} \! b_{kl} T_{ik}^{(r,m)} \overline{T}_{jl}^{(0,0)} \!\right)
\\
\phantom{S_{ij}^{(r,m)}=}{}
  + \big(q-q^{-1}\big) \sum\limits_{k,l=1}^N \sum\limits_{a+b=m-1} b_{kl} T_{ik}^{(r,a)} \overline{T}_{jl}^{(1,b)}.
\end{gather*}
\end{lemma}

\begin{proof}
Since $T_{ik}^{(0)}=0$ when $i<k$ and $\overline{T}_{jl}^{(0)}=0$ when $j>l$, we have by def\/inition
\begin{gather*}
S_{ij}^{(r,0)}= \big(q-q^{-1}\big) \sum\limits_{k,l=1}^N \sum\limits_{p=1}^{r-1} b_{kl} \left(\frac{T_{ik}^{(p)}}{q-q^{-1}}
\right) \left(\frac{\overline{T}_{jl}^{(r-p)}}{q-q^{-1}} \right)
\\
\phantom{S_{ij}^{(r,0)}=}{}
+ \big(q-q^{-1}\big)\sum\limits_{\substack{1 \leq k < i\\ 1 \leq l \leq N}} b_{kl} \left(\frac{T_{ik}^{(0)}}{q-q^{-1}} \right)
\left(\frac{\overline{T}_{jl}^{(r)}}{q-q^{-1}}\right)
\\
\phantom{S_{ij}^{(r,0)}=}{}
+ (q-1) \sum\limits_{l=1}^N b_{il} \left(\frac{T_{ii}^{(0)}-1}{q-1} \right)
\left(\frac{\overline{T}_{jl}^{(r)}}{q-q^{-1}}\right) + \sum\limits_{l=1}^N b_{il}
\frac{\overline{T}_{jl}^{(r)}}{q-q^{-1}}
\\
\phantom{S_{ij}^{(r,0)}=}{}
+ \big(q-q^{-1}\big) \sum\limits_{\substack{1 \leq k \leq N\\ j < l \leq N}} b_{kl} \left(\frac{T_{ik}^{(r)}}{q-q^{-1}} \right)
\left(\frac{\overline{T}_{jl}^{(0)}}{q-q^{-1}}\right)
\\
\phantom{S_{ij}^{(r,0)}=}{}
+ (q-1) \sum\limits_{k=1}^N b_{kj} \left(\frac{T_{ik}^{(r)}}{q-q^{-1}} \right)
\left(\frac{\overline{T}_{jj}^{(0)}-1}{q-1}\right) + \sum\limits_{k=1}^N b_{kj} \frac{T_{ik}^{(r)}}{q-q^{-1}}
\\
\phantom{S_{ij}^{(r,0)}}{}
= \sum\limits_{k=1}^N \left(b_{kj} T_{ik}^{(r,0)} + b_{ik} \overline{T}_{jk}^{(r,0)} \right) + \big(q-q^{-1}\big)
\sum\limits_{k,l=1}^N \sum\limits_{p=1}^{r-1} b_{kl} T_{ik}^{(p,0)} \overline{T}_{jl}^{(r-p,0)}
\\
\phantom{S_{ij}^{(r,0)}=}{}
+ \big(q-q^{-1}\big) \! \left(\sum\limits_{\substack{1 \leq k \leq i\\ 1 \leq l \leq N}}\! \left(\frac{q}{q+1}\right)^{\delta_{ik}}\!
b_{kl} T_{ik}^{(0,0)} \overline{T}_{jl}^{(r,0)} + \sum\limits_{\substack{1 \leq i \leq N\\ j \leq l \leq N}}\!
\left(\frac{q}{q+1}\right)^{\delta_{jl}} \! b_{kl} T_{ik}^{(r,0)} \overline{T}_{jl}^{(0,0)} \right).
\end{gather*}
This proves the case $m=0$.
The case $m=1$ is similar, except for the presence of an extra sum.
The general case follows immediately by induction on~$m$.
\end{proof}

The previous lemma along with~\eqref{TbarT} yields the next corollary.

\begin{corollary}
\label{Scong}
Under the same assumption as in the previous lemma, we have $S_{ij}^{(r,m)} \in \mathbf{K}_m$, and
\begin{gather*}
S_{ij}^{(r,m)} \equiv \sum\limits_{k=1}^N \left(b_{kj}T_{ik}^{(0,m)} + (-1)^{m+1} b_{ik} T_{jk}^{(0,m)}\right) \\
\hphantom{S_{ij}^{(r,m)} \equiv}{}
+
\big(q-q^{-1}\big) \sum\limits_{k,l=1}^N \sum\limits_{p=1}^m (-1)^{m+1-p} b_{kl} T_{ik}^{(0,p-1)} T_{jl}^{(0,m-p)}
\end{gather*}
modulo $\mathbf{K}_{m+1}$.
\end{corollary}
Again let's view $S_{ij}^{(r)}$ as an element of $\mathfrak{U}_q(\mathcal{L}(\mathfrak{gl}_N))$.
Let $\zeta_{ij}^{(r,m)}$ be the image of $S_{ij}^{(r,m)}$ in $\mathbf{K}_m/\mathbf{K}_{m+1}$.
Note that, by Corollary~\ref{Scong}, $\zeta_{ij}^{(r,m)}$ is independent of $r>0$.
\begin{theorem}
\label{twUqY}
Identify $s_{ij}^{(m+1)}$ with its image under the embedding of Proposition~{\rm \ref{embY}}.
For each $m \geq 0$, we have
\begin{gather*}
\varphi\big(s_{ij}^{(m+1)}\big)=\zeta_{ij}^{(1,m)},
\end{gather*}
where~$\varphi$ is the isomorphism of Theorem~{\rm \ref{mainthm}}.
\end{theorem}

\begin{remark}
It would also be possible to prove a~similar theorem with $\zeta_{ij}^{(0,m)}$ instead of $\zeta_{ij}^{(1,m)}$, but the
proof would be longer because we would have to consider dif\/ferent cases in order to take into account that about half of
the generators $S_{ij}^{(0)}$ are~$0$.
\end{remark}

\begin{proof}
By Proposition~\ref{embY}, Corollary~\ref{Scong} and the def\/inition of~$\varphi$, we have
\begin{gather*}
\varphi\big(s_{ij}^{(m+1)}\big)=\sum\limits_{k=1}^N \left(g_{kj} \xi_{ik}^{(0,m)} + (-1)^{m+1} g_{ik} \xi_{jk}^{(0,m)}\right)
\\
\hphantom{\varphi\big(s_{ij}^{(m+1)}\big)=}{}
+ \hbar \sum\limits_{k,l=1}^N \sum\limits_{p=1}^m (-1)^{m+1-p} g_{kl} \xi_{ik}^{(0,p-1)} \xi_{jl}^{(0,m-p)}
  = \zeta_{ij}^{(1,m)},
\end{gather*}
where we have used the fact that $b_{kl}-g_{kl}$ is divisible by $q-1$.
\end{proof}

We can obtain an analogue of Theorem~\ref{mainthm}.
From Theorem~\ref{qlim} and Proposition~\ref{embU}, we can deduce that the enveloping algebra of
$\mathfrak{g}_N^{tw}[s]$ is the limit when $q\rightarrow 1$ of $\mathfrak{U}_q(\mathcal{L}^{tw}(\mathfrak{g}_N))$ in the
following sense: if we let $\mathfrak{U}_{\mathcal{A}}(\mathcal{L}^{tw}(\mathfrak{g}_N))$ be the
$\mathcal{A}$-subalgebra of $\mathfrak{U}_q(\mathcal{L}^{tw}(\mathfrak{g}_N))$ generated by the $S_{ij}^{(r,0)}$ for all
$i$, $j$ and all $r\ge 0$, then
$\mathfrak{U}_{\mathcal{A}}(\mathcal{L}^{tw}(\mathfrak{g}_N))/(q-1)\mathfrak{U}_{\mathcal{A}}(\mathcal{L}^{tw}(\mathfrak{g}_N))$
is isomorphic to $\mathfrak{U}(\mathfrak{g}_N^{tw}[s,s^{-1}])$ (see the proof of Corollaries~3.5 and~3.12
in~\cite{MRS}).

We can def\/ine an algebra $\widetilde{Y}^{tw}(\mathfrak{g}_N)$ similarly to how we def\/ined
$\widetilde{Y}(\mathfrak{gl}_N)$.
For $m\ge 0$, in the orthogonal case, denote by $\mathsf{K}_m^{tw}$ the Lie ideal of $\mathfrak{o}_N^{tw}[s,s^{-1}]$
spanned~by
\begin{gather*}
E_{ij} s^r(s-1)^m - E_{ji} s^{-r} \big(s^{-1}-1\big)^m
\end{gather*}
for all $r \in \mathbb{Z}$.
In the symplectic case, let $\mathsf{K}_m^{tw}$ be the Lie ideal of $\mathfrak{sp}_N^{tw}[s,s^{-1}]$ spanned~by
\begin{gather*}
E_{ij'}s^r(s-1)^m - (-1)^{i+j+1} E_{ji'} s^{-r}\big(s^{-1}-1\big)^m
\end{gather*}
for all $r \in \mathbb{Z}$.
Let $U^{tw}$ be the subspace of $\mathfrak{U}_{\mathcal{A}}(\mathcal{L}^{tw}(\mathfrak{g}_N))$ spanned over $\mathbb{C}$
by all the genera\-tors~$S_{ij}^{(r,0)}$, and observe that $U^{tw}\cap
(q-1)\mathfrak{U}_{\mathcal{A}}(\mathcal{L}^{tw}(\mathfrak{g}_N)) = \{0 \}$.
Let $\mathbb{K}_m^{tw}$ be the two-sided ideal of $\mathfrak{U}_{\mathcal{A}}(\mathcal{L}^{tw}(\mathfrak{g}_N))$
generated by $\psi^{-1}(\mathsf{K}_m^{tw})\cap U^{tw}$, where~$\psi$ this time denotes the composite
\begin{gather*}
\mathfrak{U}_{\mathcal{A}}(\mathcal{L}^{tw}(\mathfrak{g}_N)) \twoheadrightarrow
\mathfrak{U}_{\mathcal{A}}(\mathcal{L}^{tw}(\mathfrak{g}_N))/(q-1)\mathfrak{U}_{\mathcal{A}}(\mathcal{L}^{tw}(\mathfrak{g}_N))
\stackrel{\sim}{\longrightarrow} \mathfrak{U}\big(\mathfrak{g}_N^{tw}\big[s,s^{-1}\big]\big).
\end{gather*}
Set $\mathbf{K}_m^{tw}$ equal to the sum of the ideals $(q-q^{-1})^{m_0} \mathbb{K}_{m_1}^{tw}\cdots
\mathbb{K}_{m_k}^{tw}$ with $m_0+m_1+\cdots + m_k \geq m$.

Let $\widetilde{Y}^{tw}(\mathfrak{g}_N)$ be the $\mathbb{C}$-algebra
\begin{gather*}
\bigoplus_{m=0}^{\infty} \mathbf{K}_m^{tw}/\mathbf{K}_{m+1}^{tw},
\end{gather*}
where $\mathbf{K}_0^{tw} = \mathfrak{U}_{\mathcal{A}}(\mathcal{L}^{tw}(\mathfrak{g}_N))$.
We also view $\widetilde{Y}^{tw}(\mathfrak{g}_N)$ as an algebra over $\mathbb{C}[\hbar]$ by setting $\hbar =
\overline{q-q^{-1}} \in \mathbf{K}_1^{tw}/\mathbf{K}_2^{tw}$.

\begin{theorem}
\label{twUqY2}
$Y^{tw}(\mathfrak{g}_N)$ is isomorphic to $\widetilde{Y}^{tw}(\mathfrak{g}_N)$ via the function $\varphi^{tw}$ that
sends $s_{ij}^{(m+1)}$ to $\overline{S_{ij}^{(1,m)}} \in \mathbf{K}_{m}^{tw}/\mathbf{K}_{m+1}^{tw}$ for $m \ge 0$.
\end{theorem}

\begin{proof}
Theorem~\ref{twUqY} implies that the following diagram is commutative:
\begin{gather*}
\xymatrix{Y^{tw}(\mathfrak{g}_N) \ar[r] \ar[d]_{\varphi^{tw}} & Y(\mathfrak{gl}_N) \ar[d]^{\varphi}\\
 \widetilde{Y}^{tw}(\mathfrak{g}_N) \ar[r] & \widetilde{Y}(\mathfrak{gl}_N).}
\end{gather*}
In this diagram, the top horizontal arrow is the embedding of Proposition~\ref{embY} and the bottom horizontal arrow is
the one induced from the embedding of Proposition~\ref{embU}.
The injectivity of~$\varphi^{tw}$ now follows from the fact that~$\varphi$ provides an isomorphism between
$Y(\mathfrak{gl}_N)$ and $\widetilde{Y}(\mathfrak{gl}_N)$: see Theorem~\ref{mainthm}.

We need to see that $\varphi^{tw}$ is surjective.
Def\/ine elements $\widetilde{S}_{ij}^{(r,m)}$ with $0 \leq r \leq m$ as follows: for each $m \geq 0$, let
$\widetilde{S}_{ij}^{(0,m)}=S_{ji}^{(0,m)}$ and $\widetilde{S}_{ij}^{(m,m)}=(-1)^{m+1}S_{ij}^{(0,m)}$, and for $1 \leq r
\leq m$ let
\begin{gather*}
\widetilde{S}_{ij}^{(r,m+1)}=\widetilde{S}_{ij}^{(r-1,m)}-\widetilde{S}_{ij}^{(r,m)}.
\end{gather*}

Then induction on~$m$ shows that
\begin{gather*}
\psi(S_{ij}^{(r,m)})=E_{ij}s^r(s-1)^m - E_{ji}s^{-r}\big(s^{-1}-1\big)^m,\\
\psi(S_{ji}^{(r,m)})=(-1)^{m+1}\big(E_{ij}s^{-(m+r)}(s-1)^m - E_{ji}s^{m+r}\big(s^{-1}-1\big)^m\big),\\
\psi(\widetilde{S}_{ij}^{(r,m)})=(-1)^{m+1}\big(E_{ij}s^{-(m-r)}(s-1)^m - E_{ji}s^{m-r}\big(s^{-1}-1\big)^m\big).
\end{gather*}
in the orthogonal case, and similarly in the symplectic case.
It follows that for any f\/ixed~$m$, the images of these elements under~$\psi$ span $\mathsf{K}_m^{tw}$, and they are all
in $U^{tw}$ by def\/inition.
Now note that for any element $X \in \psi^{-1}(\mathsf{K}_m^{tw}) \cap U^{tw}$, there is some element~$Y$ in
 \begin{gather*}
 \operatorname{span}_{\mathbb{C}} \big\{S_{ij}^{(r,m)},\widetilde{S}_{ij}^{(r,m)} \,|\, i,j=1,\ldots,N,\; r \in \mathbb{Z}_{\ge 0}\big\}
 \end{gather*}
for which $X-Y \in (q-1)\mathfrak{U}_{\mathcal{A}}(\mathcal{L}^{tw}(\mathfrak{g}_N))$.
This follows from the fact that
\begin{gather*}
\mathfrak{U}_{\mathcal{A}}\big(\mathcal{L}^{tw}(\mathfrak{g}_N)\big)/(q-1)\mathfrak{U}_{\mathcal{A}}\big(\mathcal{L}^{tw}(\mathfrak{g}_N)\big)
\stackrel{\sim}{\longrightarrow} \mathfrak{U}\big(\mathfrak{g}_N^{tw}\big[s,s^{-1}\big]\big)
\end{gather*}
is an isomorphism.
Since $X-Y$ is also in $U^{tw}$ and since $U^{tw} \cap
(q-1)\mathfrak{U}_{\mathcal{A}}(\mathcal{L}^{tw}(\mathfrak{g}_N))=\{0\}$, we see that $X=Y$.
Therefore, any element of $\mathbf{K}_m^{tw}$ is a~sum of monomials $f(q)(q-q^{-1})^{m_0}\mathcal{M}$, where $f(q) \in
\mathcal{A}$ is not divisible by $q-1$ and $\mathcal{M}=\sigma_{i_1 j_1}^{(r_1,m_1)} \cdots \sigma_{i_k j_k}^{(r_k,m_k)}$ with
\begin{gather*}
\sigma_{i_d j_d}^{(r_d,m_d)} \in \big\{S_{i_d j_d}^{(r_d,m_d)},\widetilde{S}_{i_d j_d}^{(r_d,m_d)}\big\}
\end{gather*}
and $m_0+\cdots+m_k \geq m$.
Following the same argument as in the $\mathfrak{gl}_N$ case, such a~monomial is congruent modulo
$\mathbf{K}_{m+1}^{tw}$ to
\begin{gather*}
f(1)\big(q-q^{-1}\big)^{m_0}S_{i_1 j_1}^{(1,m_1)} \cdots S_{i_k j_k}^{(1,m_k)}
\end{gather*}
up to a~sign.
The image modulo $\mathbf{K}_{m+1}^{tw}$ of this element is
\begin{gather*}
f(1)\hbar^{m_0}\overline{S_{i_1 j_1}^{(1,m_1)}} \cdots \overline{S_{i_k j_k}^{(1,m_k)}}
\end{gather*}
and this is in the image of $\varphi^{tw}$, which proves that $\varphi^{tw}$ is surjective.
\end{proof}
\begin{thebibliography}{99}
\bibitem[13]{Dr1}
Drinfeld V.G., Quantum groups, in Proceedings of the {I}nternational {C}ongress
  of {M}athematicians, {V}ols.~1,~2 ({B}erkeley, {C}alif., 1986), Amer. Math.
  Soc., Providence, RI, 1987, 798--820.

\bibitem[18]{GuMa}
Guay N., Ma X., From quantum loop algebras to {Y}angians, \href{http://dx.doi.org/10.1112/jlms/jds021}{\textit{J.~Lond.
  Math. Soc.}} \textbf{86} (2012), 683--700.

\bibitem[29]{MRS}
Molev A., Ragoucy E., Sorba P., Coideal subalgebras in quantum af\/f\/ine algebras,
  \href{http://dx.doi.org/10.1142/S0129055X03001813}{\textit{Rev. Math. Phys.}} \textbf{15} (2003), 789--822,
  \href{http://arxiv.org/abs/math.QA/0208140}{math.QA/0208140}.

\bibitem[31]{Ol}
Olshanskii G.I., Twisted {Y}angians and inf\/inite-dimensional classical {L}ie
  algebras, in Quantum Groups ({L}eningrad, 1990), \href{http://dx.doi.org/10.1007/BFb0101183}{\textit{Lecture Notes in
  Math.}}, Vol.~1510, Springer, Berlin, 1992, 104--119.

\bibitem[32]{FRT}
Reshetikhin N.Yu., Takhtadzhyan L.A., Faddeev L.D., Quantization of {L}ie groups
  and {L}ie algebras, \textit{Leningrad Math.~J.} \textbf{1} (1990), 193--225.

\end{thebibliography}
\end{document}
